\input{../doc-class-cours.tex}

\begin{document}

\begin{center}
  \textit{Le plus grand ennemi de la connaissance n'est pas l'ignorance. C'est l'illusion de la connaissance.} 
  
  \textbf{Stephen Hawking}
\end{center}

\subsection*{Ex1 - Probabilité}

\begin{multicols}{2}
On a un lance un dé à 11 faces. On observe la face du dessus. Le dé est composé : 
\begin{itemize}[label={$\bullet$}]
  \item 5 faces 1.
  \item 3 faces 2.
  \item 2 faces 5.
  \item 1 face 6.
\end{itemize} \columnbreak

\begin{enumerate}
  \item[1.] Quel est l'univers ?
  \item[2.] Quelle est la taille de l'univers ?
  \item[3.] Quelle est la probabilité de faire 1 ? 
  \item[4.] Soit $E$ l'évènement : "obtenir une face plus grande que 4". Caculer $p(E)$.  
\end{enumerate}
\end{multicols}

\subsection*{Ex2 - Ensembles}

\begin{itemize}[label={$\bullet$}]
  \item Soit $A$ l'ensemble des entiers compris entre 4 et 8. 
  \item Soit $B$ l'ensemble des entiers compris entre 7 et 10. 
\end{itemize}

\begin{multicols}{2}
\begin{enumerate}
  \item[1.] Écrire $A$. 
  \item[2.] Écrire $B$. \columnbreak
  \item[3.] Écrire $A \cap B$. 
  \item[4.] Écrire $A \cup B$. 
\end{enumerate}
\end{multicols}

\subsection*{PB1 - Urne}

\begin{multicols}{2}

On pioche une balle de couleur dans une urne. L'urne est composée : 

\begin{itemize}[label={$\bullet$}]
  \item 4 balles Rouges.
  \item 3 balles Vertes.
  \item 3 balles Bleues.
\end{itemize} \columnbreak

\begin{enumerate}
  \item[1.] Quel est l'univers ?
  \item[2.] Quelle est la probabilité de tirer une balle rouge ? 
  \item[3.] On tire une première balle bleue. On ne la remet pas dans l'urne. Quelle est la probabilité de tirer une deuxième balle balle bleue ?  
\end{enumerate}
\end{multicols}

\subsection*{PB2 - Dé inconnu}

On a un dé avec 14 faces numérotées 1,2,3 ou 4. On connaît quelques probabilités : 

\begin{multicols}{3}
\begin{itemize}[label={$\bullet$}]
  \item $p(1) = \dfrac{1}{2}$
  \item $p(2) = \dfrac{1}{7}$
  \item $p(3) = \dfrac{1}{14}$
\end{itemize}
\end{multicols}

\begin{enumerate}
  \item[1.] Calculer $p(4)$.
  \item[2.] En déduire le nombre de faces numérotées 4.
\end{enumerate}


\subsection*{PB3 - podium}
Les trois figures sont des carrés. Calculer l'aire total du podium gris. \\

\begin{figure}[H]
    \centering
    \includegraphics[width=0.5\linewidth]{4x2-probabilites/pb3.pdf}
\end{figure}

\end{document}